% !TeX program = xelatex
% ============================================================================
%  第 9 课　D₄ 的子群与二面体群的关系　—— 学生版讲义
%  与 02-课件/第09课-D4的子群/第09课-D4的子群.tex 同步
% ============================================================================

\xhlesson{9}{\texorpdfstring{\(D_4\)}{D4} 的子群与二面体群的关系}{}

\begin{mubiao}
  \begin{itemize}
    \item 能找出 \(D_4\) 里面能自成一体的小群（子群）。
    \item 能说出旋转与翻折之间那条关系：\(fr=r^{-1}f\)。
    \item 知道 \(D_n\) 叫二面体群，正 \(n\) 边形有 \(2n\) 个对称动作。
  \end{itemize}
\end{mubiao}

\section*{一、找 \(D_4\) 的子群}

\noindent 办法照旧：圈出一小堆动作，把复合填成一张表。全在里面，就是子群。

\begin{dongshou}[先检查两个]
  \begin{enumerate}
    \item \(H=\{e,\ r^2\}\)：
      \(r^2\cdot r^2=\) \underline{\hspace{2em}}，跑出去了吗？\underline{\hspace{3em}}
      \quad 所以它\underline{\hspace{2em}}（是／不是）子群。
    \item \(H=\{e,\ r\}\)：
      \(r\cdot r=\) \underline{\hspace{2em}}，这个结果在 \(H\) 里吗？\underline{\hspace{3em}}
      \quad 所以它\underline{\hspace{2em}}（是／不是）子群。
  \end{enumerate}
\end{dongshou}

\begin{example}[\(D_4\) 里全部的子群]
  \begin{center}
  \begin{tabular}{@{}clc@{}}
    \toprule
    大小 & 子群 & 个数 \\
    \midrule
    \(1\) & \(\{e\}\) & \(1\) \\
    \(2\) & \(\{e,r^2\}\)；\(\{e,f\}\)；\(\{e,rf\}\)；\(\{e,r^2f\}\)；\(\{e,r^3f\}\) & \(5\) \\
    \(4\) & \(\{e,r,r^2,r^3\}\)；\(\{e,r^2,f,r^2f\}\)；\(\{e,r^2,rf,r^3f\}\) & \(3\) \\
    \(8\) & \(D_4\)（全部八个） & \(1\) \\
    \bottomrule
  \end{tabular}
  \end{center}
  一共 \(10\) 个。
\end{example}

\begin{sikao}
  看最大子群和最小子群的大小：\(1,2,4,8\)。它们和 \(D_4\) 的大小 \(8\) 有什么关系？\\[4pt]
  （和上节课 \(S_3\) 里的发现比一比。）
  \liukong{3}
\end{sikao}

\section*{二、旋转与翻折的关系}

\begin{dongshou}[左右各走一条路]
  拿纸片做两遍，把终点记下来：
  \begin{itemize}
    \item 先翻 \(f\)，再转 \(90^\circ\)（\(r\)）：终点在 \underline{\hspace{3em}}；
    \item 先转 \(270^\circ\)（\(r^3\)），再翻 \(f\)：终点在 \underline{\hspace{3em}}。
  \end{itemize}
  两次的终点一样吗？\underline{\hspace{3em}}
\end{dongshou}

\begin{example}
  \[
    fr=r^{-1}f, \qquad\text{也就是}\qquad fr=r^3f .
  \]
  白话：「先翻后转，等于先\textbf{倒着}转、再翻。」
\end{example}

\begin{dingli}[二面体群]
  正 \(n\) 边形的全部对称动作，按复合构成一个群，叫做\textbf{二面体群}，记作 \(D_n\)，
  一共有 \(2n\) 个动作：\(n\) 个转动加 \(n\) 个翻折。
\end{dingli}

\begin{sikao}
  \(D_3\) 有多少个动作？\(D_5\)（正五边形）呢？\(D_{12}\) 呢？
  \liukong{2}
\end{sikao}

\begin{xiaojie}
  用一句话回答：为什么说「\(D_4\) 的 \(8\) 个动作，两个就够生成了」？
  \liukong{3}
\end{xiaojie}

\noindent\textbf{下节课}：生成元、循环群与循环群的子群——为什么两个就够。